\begin{longtable}{@{}l p{12.5cm}@{}}
  \caption{Skills per-instance answer comparison ($25\% \leq \Delta < 50\%$) for the 7 most divergent items: same question, each model's answer and reasoning.}\label{tab:cmp_skills_pos25} \\
  \toprule
  \endfirsthead
  \toprule
  \endhead
  \multicolumn{2}{@{}l}{\textbf{116190} \quad $\Delta = +40\%$} \\
  \textbf{Question} & You are a research assistant given a curated list of recent peer-reviewed articles (published within the last five years) related to K-12 physical education aligned with SHAPE America National PE Standards. Using the provided article list, produce a structured literature-review synthesis that: (a) applies the SHAPE America standards to categorize the studies' primary focus areas, (b) analyzes patterns and gaps across the literature with regard to content areas (e.g., motor skills, personal and social responsibility, fitness, cognitive skills), populations (grade levels, socio-economic status, special needs), and research methods, and (c) proposes two concrete, novel research questions addressing identified gaps that could be pursued in a future empirical study. In your response, include: a brief summary (2--3 sentences) of how you mapped articles to specific SHAPE standards; a comparative analysis (3--5 paragraphs) highlighting major trends and three clear research gaps; and two well-justified research questions with a one-paragraph rationale for each including suggested population and method. Assume the article list includes mixed-methods classroom interventions, large-scale survey studies of PE curriculum fidelity, and small qualitative studies of teacher perspectives. \\
  \textbf{Reference} & Mapping to SHAPE Standards: I mapped each article to the most relevant SHAPE America National PE Standard by extracting each study's stated learning objective or measured outcomes and aligning those with the standards' domains (e.g., Standard 1: motor skills and movement patterns; Standard 2: knowledge of movement concepts; Standard 3: physical activity and fitness; Standard 4: personal and social responsibility; Standard 5: value of physical activity). Classroom intervention studies mostly aligned with Standard 1 and 3 (skill development and fitness outcomes), curriculum-fidelity surveys typically addressed Standard 2 and 5 (knowledge and valuing activity), and teacher-perspective qualitative studies emphasized Standard 4 (personal/social behaviors and classroom climate). This mapping allowed cross-study comparison by standard focus rather than by nominal topic labels used by authors. \newline Comparative analysis and trends: First, a prominent trend is the dominance of interventions targeting fundamental motor skills and fitness outcomes in elementary grades. Mixed-method classroom interventions reported measurable gains in motor skill performance and short-term fitness improvements, often using pre-post designs with objective skill rubrics or fitness tests. This emphasis indicates a strong applied interest in Standard 1 and Standard 3 outcomes for younger students. Second, large-scale survey studies on curriculum fidelity revealed moderate-to-low adherence to SHAPE-recommended curricula at the middle and high school levels, with common barriers cited as time constraints, inadequate equipment, and competing academic priorities. These surveys tended to measure teacher self-report of content delivered (linking to Standard 2 and 5) but rarely tied fidelity measures to student learning outcomes. Third, qualitative studies of teacher perspectives highlighted social and behavioral domains---classroom management, fostering cooperation, and inclusion strategies---consistent with Standard 4, but these studies were small and localized, limiting generalizability. Methodologically, there is a split: robust quantitative measures for physical outcomes, less rigorous or less frequent longitudinal follow-up, and a smaller but rich set of qualitative insights into classroom processes. \newline Across these trends, three clear gaps emerge. Gap 1: Integration across standards --- Few studies examine multi-standard interventions that simultaneously target motor competence, cognitive understanding of movement concepts, lifetime physical activity attitudes, and social responsibility; most focus on a single standard/outcome. Gap 2: Equity and diverse populations --- The literature underrepresents rigorous studies of culturally responsive PE approaches and of students with disabilities or from low-SES backgrounds; when included, these groups are often subsamples without tailored intervention designs. Gap 3: Longitudinal and fidelity-linked outcomes --- Curriculum-fidelity surveys rarely connect fidelity to longitudinal student outcomes, and intervention studies seldom include long-term follow-up beyond one semester, limiting evidence about sustained behavior change or retention of skills. \newline Research question 1 (addressing Gap 1 and 3): "Does a multi-component PE curriculum explicitly designed to integrate motor skill instruction (Standard 1), movement concepts (Standard 2), and social responsibility activities (Standard 4) produce greater improvements in student motor competence, movement knowledge, and prosocial behaviors over one academic year compared to a standard single-focus PE curriculum?" Rationale: This question tests whether integrating multiple SHAPE standards yields additive or synergistic benefits versus the common single-outcome approach. Suggested population and method: a cluster randomized controlled trial in upper elementary schools (grades 4--5), with schools randomized to the integrated curriculum or business-as-usual PE. Measures would include validated motor competence assessments, a standardized movement-concepts test, teacher- and observer-rated social-responsibility behavior scales, and pre-post measures plus a six-month follow-up to assess retention. Fidelity monitoring protocols would document implementation quality to permit fidelity-outcome analyses. \newline Research question 2 (addressing Gap 2): "How does a culturally adapted PE intervention designed for low-SES, racially diverse middle-school students, including differentiated instruction for students with mild-to-moderate disabilities, affect students' activity engagement, attitudes toward lifelong physical activity (Standard 5), and equitable access to skill development opportunities?" Rationale: This question targets equity by evaluating whether culturally responsive adaptations and differentiated instruction increase engagement and perceived value of physical activity among historically underserved groups. Suggested population and method: a mixed-methods quasi-experimental design in urban middle schools, with intervention and matched comparison schools. Quantitative measures would include device-based activity monitoring during PE, attitude and value surveys, and skill opportunity indices (e.g., proportion of class time each student engaged in skill practice). Complementary qualitative data from focus groups with students and teachers would explore perceived relevance and barriers, informing scalability and refinement. Together, these two questions address major gaps by testing integrated, standards-aligned curricula and by centering equity and sustained outcomes in K--12 PE research. \\
  \midrule
  \textbf{Kimi-K2.6} & score 5/10 \\
  & The user wants a structured literature-review synthesis based on a curated list of recent peer-reviewed articles related to K-12 physical education aligned with SHAPE America National PE Standards. However, the user hasn't actually provided the specific article list. They said "Using the provided article list" but no list was provided in the prompt. \newline I need to handle this carefully. I have two options: \newline 1. Ask the user to provide the list \newline 2. Create a hypothetical but realistic synthesis based on the description they gave (mixed-methods classroom interventions, large-scale survey studies of PE curriculum fidelity, small qualitative studies of teacher perspectives) and the SHAPE America standards. \newline Given the instructions say "Using the provided article list" but none was provided, and the user says "Assume the article list includes..." at the end, this suggests I should construct a plausible synthesis based on typical recent literature in the field and the methodological mix described. The "Assume" clause gives me permission to work with a hypothetical corpus matching that description. \newline SHAPE America National Standards for K-12 Physical Education (revised typically): \newline \textit{[\ldots\ truncated \ldots]} \newline 3. Proposed Research Questions (two sections with RQ and rationale paragraph) \newline Looks good. I'll write it out now. \\
  \textbf{Innospark-1T} & score 9/10 \\
  & Let me think through this carefully. \newline The question/task: You are a research assistant given a curated list of recent peer-reviewed articles (published within the last five years) related to K-12 physical education aligned with SHAPE America National PE Standards. Using the provided article list, produce a structured literature-review synthesis that: (a) applies the SHAPE America standards to categorize the studie \newline My reasoning: Literature Review Synthesis: K-12 Physical Education and SHAPE America Standards \newline Mapping Articles to SHAPE America Standards \newline I mapped the 24 peer-reviewed articles (2019--2024) to the five SHAPE America National Physical Education Standards by analyzing each study's primary outcomes, interventions, and measured variables. Standard 5 (valuing physical activity) was least represented, with only 3 articles primarily addressing this domain. Research Gaps Identified \newline 1. Limited Longitudinal Research: Only 2 studies followed participants beyond 12 weeks, making it impossible to assess sustained behavior change or skill retention. Standard 5 (valuing physical activity) was least represented, with only 3 articles primarily addressing this domain. Research Gaps Identified \newline 1. \\
  \specialrule{1.2pt}{8pt}{8pt}
  \multicolumn{2}{@{}l}{\textbf{113936} \quad $\Delta = +40\%$} \\
  \textbf{Question} & You are a research assistant compiling an accessible list of K-12 music education research ideas aligned with the National Core Arts Standards. Using recent literature trends, analyze gaps in current research and propose five distinct, researchable study ideas that an educator could realistically implement. For each proposed idea, break down (a) the specific research question, (b) the rationale grounded in identified gaps or trends, (c) the target population and setting, and (d) the primary research method(s) and possible data sources. Ensure your analysis highlights relationships among the ideas (e.g., how one study could inform another) and explains why each idea is both feasible for teachers and likely to contribute new knowledge to K-12 music education. \\
  \textbf{Reference} & Model answer: \newline Analysis of current trends and gaps: Recent literature in K-12 music education shows growing interest in technology-enhanced instruction, culturally responsive pedagogy, social-emotional learning (SEL) through music, inclusive practices for students with disabilities, and assessment alternatives (e.g., formative, performance-based assessments). However, gaps remain in (1) classroom-scalable evaluations of specific tech tools across grade bands, (2) empirical comparisons between culturally responsive curricula and standard curricula on engagement and achievement, (3) practical models linking music-making to SEL outcomes with validated measures, (4) classroom-level strategies for differentiating music instruction for neurodiverse learners, and (5) teacher-built, classroom-feasible assessment approaches that align with National Core Arts Standards and are shown to impact instruction. \newline Based on these gaps, I propose five researchable study ideas: \newline 1) Study idea 1 --- Research question: Does integrating tablet-based music creation apps into elementary general music lessons improve rhythmic accuracy and creative expression compared with traditional notation-based activities? \newline (a) Rationale: Many studies examine technology broadly but few evaluate specific apps in tangible musical outcomes for early grades. This fills the gap of classroom-scalable tech efficacy. \newline (b) Target population/setting: Grade 3--4 general music classes in 4--6 public elementary schools, mixed socioeconomic contexts. \newline \textit{[\ldots\ truncated \ldots]} \newline (d) Relationship to others: This protocol can be used as the assessment mechanism in the other four studies to provide consistent outcome measures and to study how assessment-informed instruction mediates student gains. \newline Synthesis and feasibility considerations: These five studies target complementary gaps---technology efficacy, cultural relevance, SEL impacts, inclusive differentiation, and formative assessment. They are designed for classroom feasibility: short units, teacher-led interventions, common-grade settings, and practical measures (recordings, rubrics, surveys, observations). Together, they create a coherent research agenda where outcomes from one study (e.g., effective differentiation or assessment protocols) can be leveraged in others, creating cumulative knowledge that directly supports educators implementing the National Core Arts Standards. \\
  \midrule
  \textbf{Kimi-K2.6} & score 5/10 \\
  & The user wants a research assistant response that compiles accessible K-12 music education research ideas aligned with the National Core Arts Standards (NCAS). Requirements: \newline 1. Use recent literature trends. \newline 2. Analyze gaps in current research. \newline 3. Propose five distinct, researchable study ideas that an educator could realistically implement. \newline 4. For each idea, break down: \newline - (a) specific research question \newline \textit{[\ldots\ truncated \ldots]} \newline Need to ensure this is concise. Maybe use a table? The user asked to break down a,b,c,d for each. A structured list is better than a table for readability. \newline Let's write it out cleanly. \\
  \textbf{Innospark-1T} & score 9/10 \\
  & Let me think through this carefully. \newline The question/task: You are a research assistant compiling an accessible list of K-12 music education research ideas aligned with the National Core Arts Standards. Using recent literature trends, analyze gaps in current research and propose five distinct, researchable study ideas that an educator could realistically implement. For each proposed idea, break down (a) the spe \newline My reasoning: Research Ideas in K-12 Music Education \newline Analysis of Current Literature Trends and Gaps \newline Recent Trends in Music Education Research: \newline 1. Technology integration in music learning \newline \textit{[\ldots\ truncated \ldots]} \newline Conclusion \newline These five research ideas address significant gaps in current literature while remaining feasible for practicing educators. Each aligns with specific National Core Arts Standards anchor standards (Creating, Performing, Responding, Connecting) and employs methodologies appropriate for classroom-based inquiry. The interconnected nature of these studies allows educators to build a program of research that contributes meaningfully to both local practice and the broader field of music education. \\
  \specialrule{1.2pt}{8pt}{8pt}
  \multicolumn{2}{@{}l}{\textbf{109515} \quad $\Delta = +40\%$} \\
  \textbf{Question} & You are evaluating a learner's short essay response to a prompt about reaction mechanisms. The learner wrote the following three-sentence answer about the mechanism of nucleophilic acyl substitution of an ester by hydroxide in basic hydrolysis: \newline "1. The hydroxide ion attacks the carbonyl carbon, forming a tetrahedral intermediate. 2. The tetrahedral intermediate then eliminates methoxide ion directly, regenerating the carbonyl and producing methanol. 3. Because the reaction occurs in base, the final product is a neutral carboxylic acid." \newline Analyze each sentence for factual accuracy and mechanistic errors. For each sentence, (a) identify any mistake(s) or incomplete reasoning, (b) explain the correct mechanism or concept that should replace or supplement the learner's statement, and (c) indicate how the error would affect the predicted products or reaction pathway if left uncorrected. Provide a concise, coherent corrected version of the entire three-sentence answer that is mechanistically accurate and appropriate for a higher-education chemistry audience. \\
  \textbf{Reference} & Sentence 1 analysis: \newline (a) Mistake/incompleteness: The statement that "hydroxide ion attacks the carbonyl carbon, forming a tetrahedral intermediate" is essentially correct but incomplete. It omits that the attack increases electron density on the carbonyl oxygen, and that the intermediate is negatively charged (alkoxide-like). It also doesn't state that this step is reversible and that the reaction coordinate includes proton transfers in the solvent. \newline (b) Correct mechanism/concept: Hydroxide (a strong nucleophile) adds to the electrophilic carbonyl carbon of the ester, pushing the C=O electrons onto oxygen to give a tetrahedral alkoxide intermediate. The intermediate bears a negative charge on oxygen and is stabilized by solvation; the formation of this intermediate is generally reversible under basic conditions. \newline (c) Effect if uncorrected: Omitting reversibility and the alkoxide character could lead to misunderstandings about subsequent steps (e.g., the ease of collapse of the intermediate) and the role of solvent and proton transfer in determining product distribution. \newline Sentence 2 analysis: \newline (a) Mistake/incompleteness: The learner states that the tetrahedral intermediate "then eliminates methoxide ion directly, regenerating the carbonyl and producing methanol." This is incorrect in two ways: (1) In basic hydrolysis, the leaving group is expelled as an alkoxide (methoxide), not as methanol; methanol would only form if methoxide were protonated. (2) The expulsion of methoxide is driven by reformation of the carbonyl, but because the medium is basic, methoxide remains deprotonated and does not convert to methanol without an acid workup. \newline \textit{[\ldots\ truncated \ldots]} \newline Coherent corrected three-sentence answer: \newline 1. Hydroxide ion attacks the ester carbonyl carbon, forming a negatively charged tetrahedral (alkoxide) intermediate in a generally reversible addition step. 2. The tetrahedral intermediate collapses to re-form the carbonyl and expels methoxide ion (CH3O) as the leaving group; in the basic medium the expelled species remains deprotonated. 3. Consequently, the immediate organic product is the carboxylate anion (not the neutral carboxylic acid); neutral acid is obtained only after subsequent acid workup. \\
  \midrule
  \textbf{Kimi-K2.6} & score 5/10 \\
  & [5-6] Identifies major errors but introduces a mechanistic inaccuracy in Sentence 2's correction, claiming the intermediate collapses to a carboxylic acid which is then deprotonated to yield methanol. The reference specifies the collapse directly expels methoxide, yielding carboxylate, not methanol. Sentence 1 analysis also lacks depth on reversibility. \\
  \textbf{Innospark-1T} & score 9/10 \\
  & Let me think through this carefully. \newline The question/task: You are evaluating a learner's short essay response to a prompt about reaction mechanisms. The learner wrote the following three-sentence answer about the mechanism of nucleophilic acyl substitution of an ester by hydroxide in basic hydrolysis: \newline "1. The hydroxide ion attacks the carbonyl carbon, forming a tetrahedral intermediate. \newline 2. The tetrahedral intermediate then eliminates methoxide ion directly, regenerating the carbonyl and producing methanol. 3. Because the reaction occurs in base, the final product is a neutral carboxylic acid." \newline My reasoning: Analysis of the Learner's Response on Ester Hydrolysis Mechanism \newline Sentence 1: "The hydroxide ion attacks the carbonyl carbon, forming a tetrahedral intermediate." \newline \textit{[\ldots\ truncated \ldots]} \newline This error demonstrates a fundamental misunderstanding of acid-base chemistry in reaction contexts. Corrected Version: \newline "1. The hydroxide ion attacks the electrophilic carbonyl carbon, pushing electrons onto the oxygen to form a tetrahedral alkoxide intermediate. 2. The tetrahedral intermediate then collapses, with the carbonyl oxygen's lone pair reforming the C=O double bond while simultaneously expelling the methoxide leaving group, producing a carboxylic acid and methoxide. 3. \\
  \specialrule{1.2pt}{8pt}{8pt}
  \multicolumn{2}{@{}l}{\textbf{13748} \quad $\Delta = +40\%$} \\
  \textbf{Question} & You are an instructional designer tasked with writing a needs analysis report for a university-level philosophy program that aims to integrate practical ethics and moral formation into the curriculum. Based on interviews with stakeholders (faculty, students, employers), review of institutional documents, and secondary literature, produce a comprehensive needs analysis report section that: (1) identifies and justifies the program's core pedagogical needs by tracing root causes (use causal analysis to distinguish surface symptoms from underlying issues); (2) formulates clear program goals that reflect both cognitive and character development (internalizing moral values) and explains how these goals align with institutional mission and stakeholder expectations; (3) proposes at least three prioritized, evidence-based solutions (curriculum design changes, assessment strategies, faculty development, or co-curricular initiatives), justifying each choice with expected impacts on learning and internalization of values; and (4) anticipates potential implementation challenges and offers concrete mitigation strategies. Your response should evaluate trade-offs, demonstrate synthesis of qualitative and quantitative evidence, and articulate how the proposed solutions will foster students' sustained commitment to ethical dispositions rather than only performative compliance. \\
  \textbf{Reference} & Needs Analysis Report --- Integrating Practical Ethics and Moral Formation into the Philosophy Curriculum \newline 1. Identification and causal analysis of core pedagogical needs \newline Summary of findings: Stakeholder interviews and document review reveal three recurring problems: (a) graduates show strong analytical competence but limited demonstration of consistent ethical dispositions in professional contexts; (b) students report limited opportunities for applied, iterative practice in ethical decision-making; (c) faculty express uneven confidence in facilitating values-focused, affective learning and in assessing moral development. Quantitative indicators (alumni employer surveys: 62\% rated graduates' ethical judgment as only adequate; capstone reflection analyses: 48\% of students failed to demonstrate integrated moral commitments across cases) corroborate these themes. \newline Causal analysis (root causes vs. symptoms): \newline - Surface symptom: Poor demonstration of internalized ethical dispositions in professional settings. \newline - Root cause 1: Curriculum emphasis on abstract analytic skills with insufficient scaffolded, situated practice (course sequencing privileges theory-heavy content and limited applied seminars). Evidence: program syllabus audit shows 70\% of required credits in traditional text-based courses and only one required applied ethics seminar. \newline \textit{[\ldots\ truncated \ldots]} \newline Conclusion --- How proposed solutions foster sustained internalization \newline By reconfiguring curriculum to provide repeated, authentic practice; reforming assessment to value dispositional evidence; and building faculty capacity and co-curricular supports, the program moves from episodic instruction to an ecology of learning that cultivates moral habits. The integrated approach addresses root causes identified in the causal analysis: students will encounter ethically charged situations across time, receive formative and summative feedback on behavior and reflection, and be guided by trained faculty and mentors. Over time, such systemic alignment---measured by triangulated outcomes (portfolios, observed performance, employer reports)---should shift student motivations from external performance to internalized ethical commitments, thereby achieving the program goals of cognitive integration and character formation. \\
  \midrule
  \textbf{Kimi-K2.6} & score 5/10 \\
  & [5-6] The answer is incomplete, cutting off midway through Priority 2. While sections 1 and 2 effectively identify needs, root causes, and goals with good evidence, the missing half means only some key requirements are captured. Priorities, constraints, success indicators, and implementation challenges are insufficiently developed or absent. \\
  \textbf{Innospark-1T} & score 9/10 \\
  & Let me think through this carefully. \newline The question/task: You are an instructional designer tasked with writing a needs analysis report for a university-level philosophy program that aims to integrate practical ethics and moral formation into the curriculum. Based on interviews with stakeholders (faculty, students, employers), review of institutional documents, and secondary literature, produce a comprehensive needs analysis report section that: (1) ide \newline My reasoning: Needs Analysis Report: Integrating Practical Ethics and Moral Formation in University Philosophy Curriculum \newline 1. Root Cause: Epistemological Fragmentation \newline - Students experience ethics as abstract intellectual puzzles disconnected from lived experience \newline - Faculty disciplinary specialization creates siloed moral reasoning without integrative frameworks \newline \textit{[\ldots\ truncated \ldots]} \newline Conclusion: \newline This needs analysis reveals that effective moral formation requires not merely content coverage but pedagogical transformation---shifting from knowledge transmission to situated practice, from individual assessment to community formation, and from compliance-oriented evaluation to developmental partnership. The proposed solutions, while requiring significant institutional investment, offer the greatest probability of producing graduates whose ethical commitments are internally motivated, intellectually grounded, and practically actionable. \\
  \specialrule{1.2pt}{8pt}{8pt}
  \multicolumn{2}{@{}l}{\textbf{21508} \quad $\Delta = +30\%$} \\
  \textbf{Question} & As an instructional designer creating a K--12 mathematics lesson aligned to a specific affective learning objective written in Gagn's ABCD format (Audience, Behavior, Condition, Degree) about students demonstrating positive attitudes toward collaborative problem-solving, describe in detail: 1) a clear ABCD-stated affective learning objective for this lesson; 2) the instructional strategies and media you would select to support that objective; and 3) a step-by-step lesson plan mapped to Gagn's nine events of instruction that shows how you will engage students at the Responding level of the affective domain (active participation and reaction). Be specific about examples of activities, teacher prompts, student responses, and how each Gagn event supports the affective goal. \\
  \textbf{Reference} & 1) ABCD-stated affective learning objective: \newline Audience: 5th-grade math students in a heterogeneous classroom. \newline Behavior: demonstrate a positive attitude toward collaborative problem-solving by actively contributing ideas, listening respectfully to peers, and encouraging group members. \newline Condition: given a pair or small-group mathematics task involving a multi-step word problem and teacher-facilitated reflection prompts. \newline Degree: in at least 4 out of 5 group sessions during the unit, as observed by teacher checklists and a short self-reflection, students will display the behaviors (contribution, respectful listening, encouragement) for a minimum of 8 minutes per 20-minute group work period. \newline 2) Instructional strategies and media: \newline \textit{[\ldots\ truncated \ldots]} \newline How it supports the objective: fosters transfer of affective responses to new contexts and encourages repeated practice. \newline Overall, this plan keeps the affective learning objective central by repeatedly prompting active participation, modeling and reinforcing pro-social language, and providing both teacher and peer feedback. The structured roles and short reflection tasks reduce ambiguity about expected behaviors, while the observation checklist and self-reflection measure responding-level engagement so the teacher can adjust supports in subsequent lessons. \\
  \midrule
  \textbf{Kimi-K2.6} & score 5/10 \\
  & The user wants a detailed K-12 mathematics lesson plan with three specific components: \newline 1. An ABCD-stated affective learning objective about positive attitudes toward collaborative problem-solving. \newline 2. Instructional strategies and media to support that objective. \newline 3. A step-by-step lesson plan mapped to Gagn's nine events of instruction, engaging students at the Responding level of the affective domain (Krathwohl's taxonomy: receiving, responding, valuing, etc.). \newline Key constraints: \newline - Must be K-12 mathematics (I should pick a specific grade level to make it concrete; let's say 7th grade math - proportional reasoning/ratios, as it's collaborative-friendly). \newline \textit{[\ldots\ truncated \ldots]} \newline One final check: Is "reacting supportively" at the Responding level? Yes. Responding includes acquiescing, willing participation, and satisfaction in response. Active contribution and supportive reaction are classic Responding indicators. Valuing would be if they consistently choose to work in groups or believe collaboration is important even when not asked. Responding is the "doing" and "reacting" part, which fits perfectly. \newline I think I'm ready to write. \\
  \textbf{Innospark-1T} & score 8/10 \\
  & Let me think through this carefully. \newline The question/task: As an instructional designer creating a K--12 mathematics lesson aligned to a specific affective learning objective written in Gagn's ABCD format (Audience, Behavior, Condition, Degree) about students demonstrating positive attitudes toward collaborative problem-solving, describe in detail: 1) a clear ABCD-stated affective learning objective for this lesson; 2) the \newline My reasoning: Instructional Design: Affective Learning in Mathematics \newline 1) ABCD-Stated Affective Learning Objective \newline Objective: \newline "Given a collaborative group setting with assigned roles and a challenging multi-step mathematics problem (Condition), middle school students (Audience) will actively participate in group discussions, offer constructive feedback to peers, and demonstrate willingness to incorporate others' ideas (Behavior) by contributing at least one substantive idea or question per problem-solving session and maintaining positive body language and tone when interacting with teammates, as observed by the teacher using a collaborative rubric (Degree)." \newline \textit{[\ldots\ truncated \ldots]} \newline | Assessment | Rubric measures observable collaborative behaviors, not just attitudes | \newline This lesson design ensures students move beyond mere receiving (passive awareness) to responding (active participation) in collaborative problem-solving, laying groundwork for future valuing and organization of collaborative skills. \\
  \specialrule{1.2pt}{8pt}{8pt}
  \multicolumn{2}{@{}l}{\textbf{97341} \quad $\Delta = +30\%$} \\
  \textbf{Question} & You are an evaluator conducting a summative assessment of a kindergarten physical development \& health program. Using rubric-derived strengths and weaknesses across affective-domain criteria (e.g., cooperation, safety attitudes, persistence, respect for rules), analyze and prioritize three programmatic values that should guide revisions to the curriculum. For each prioritized value: (1) explain why it takes priority based on the rubric evidence, (2) describe how conflicts among values (for example, fostering independence vs. ensuring safety, or competition vs. cooperation) should be resolved in classroom practice, and (3) propose two concrete, affective-focused changes to the program that align with the prioritized value and address identified weaknesses. Support your analysis with specific examples drawn from rubric evidence and projected classroom scenarios. \\
  \textbf{Reference} & Based on the rubric evaluation of the kindergarten physical development \& health program, the three prioritized affective values for curriculum revision are: 1) Safety-mindedness, 2) Cooperative engagement, and 3) Resilient persistence. These priorities follow from rubric strengths and weaknesses: children scored highest on basic motor participation but lower on consistent safety behaviors and respectful turn-taking; persistence during challenging tasks showed mixed results with many children abandoning tasks when unsuccessful. Prioritizing these values will address the program's most consequential weaknesses while preserving existing strengths. \newline 1) Safety-mindedness \newline Why priority: The rubric shows recurring weaknesses in 'recognizing and following safety rules' and 'using equipment safely' across several assessment tasks (e.g., 40\% of children needed adult prompts to store equipment; 35\% attempted risky behaviors during obstacle courses). Since safety lapses can cause injury and limit participation, strengthening safety-mindedness is foundational. \newline Resolving conflicts: The main tension is between promoting independence and enforcing safety protocols. Resolve this by shifting from adult-imposed restrictions to scaffolded autonomy: teach safety rules explicitly, model safe choices, and give children graduated responsibility (e.g., choosing which safe path in an obstacle course). Use guided reflection after activities so children evaluate risks themselves, fostering internalized safety values rather than compliance out of fear. \newline Two concrete changes: \newline - Introduce a 'safety leader' rotation where small teams jointly inspect and set up equipment under teacher supervision, then reflect on risks and safe practices. This builds ownership and procedural knowledge while preserving safety oversight. \newline \textit{[\ldots\ truncated \ldots]} \newline - Use structured failure-safe activities where challenges are framed as experiments (e.g., "Let's try three ways to balance on one foot"), encouraging children to try multiple strategies and verbalize what they learned, which reduces fear of failure and builds persistence. \newline Integration and examples: For example, in an obstacle-course unit the rubric showed safe handling but poor turn-taking and quitting when tasks seemed hard. Applying the prioritized values, redesign the unit so teams collaboratively plan a safe route (safety-mindedness), rotate roles so each child leads once and supports peers (cooperative engagement), and set stepwise balancing goals with praise for attempts (resilient persistence). Conflicts---such as a child wanting to rush through a challenge to win---are managed by team norms that reward helping over speed and by teacher prompts that reframe success as everyone completing their step safely. Over time, these changes should increase internalization of safety rules, improve social interactions during physical activities, and raise persistence when learning new motor skills, directly addressing the weaknesses identified in the rubric while maintaining high levels of active participation. \\
  \midrule
  \textbf{Kimi-K2.6} & score 6/10 \\
  & [5-6] The answer partially meets the basic requirements as it is incomplete, cutting off midway through the second value and omitting the third entirely. However, the provided content is highly detailed, logically coherent, and features creative insights (e.g., 'scaffolded autonomy', 'Traffic Light Body Check-In'). The incompleteness prevents a higher score despite the strong quality of the existing sections. \\
  \textbf{Innospark-1T} & score 9/10 \\
  & Let me think through this carefully. \newline The question/task: You are an evaluator conducting a summative assessment of a kindergarten physical development \& health program. Using rubric-derived strengths and weaknesses across affective-domain criteria (e.g., cooperation, safety attitudes, persistence, respect for rules), analyze and prioritize three programmatic values that should guide revisions to the curriculum. For each prioritized value: (1) expl \newline My reasoning: Summative Assessment of Kindergarten Physical Development \& Health Program \newline Analysis of Rubric Evidence \newline Based on hypothetical rubric data from classroom observations, teacher assessments, and student performance indicators, I've identified the following patterns: \newline | Affective Domain Criteria | Strength Level | Weakness Indicators | \newline \textit{[\ldots\ truncated \ldots]} \newline - Example: Before climbing structure, students complete checklist: "Hands ready? Eyes watching? Space clear?" \newline 2. \\
  \specialrule{1.2pt}{8pt}{8pt}
  \multicolumn{2}{@{}l}{\textbf{41448} \quad $\Delta = +30\%$} \\
  \textbf{Question} & You are a teaching assistant monitoring a university-level ensemble rehearsal. During a live performance, a section of student musicians consistently plays a passage one beat late and with incorrect articulation compared to the score. Describe, in detail, how you would detect and analyze the root causes of this error in real time and provide at least three concrete corrective suggestions you would deliver immediately to the instructor to help the students fix the issue during the same rehearsal. In your analysis, break down possible sensory, cognitive, and pedagogical factors (e.g., score reading, ensemble cues, rehearsal technique, tempo perception, notation misinterpretation) and explain how each factor could produce the observed late entries and wrong articulation. For each corrective suggestion, explain why it targets specific root causes, how the instructor should implement it in the moment (exact phrasing or gesture if relevant), and what measurable cue or behavior the instructor should look for to confirm improvement within the rehearsal. \\
  \textbf{Reference} & Detection and analysis: \newline First, I would confirm the pattern by listening for consistency rather than isolated instances: the passage is one beat late and has an articulation mismatch across multiple measures and multiple players in the same section. I would quickly check the score at that passage to confirm the written articulation and placement of notes relative to the beat (e.g., is there an anacrusis, syncopation, or pickup that could be misread?). I would observe the players' body language and breath/cueing behavior---are they watching the conductor, a section leader, or the prior barline? Are they anticipating a cue or reacting late? I would also note the tempo and conductor gestures: is the tempo steady, or has the conductor's preparatory beat been ambiguous? Finally, I would consider rehearsal history: was that passage recently changed (dynamic/articulation) or has it been rehearsed less than surrounding material? \newline Possible root causes (broken down): \newline - Score-reading/notation misinterpretation: The passage might include syncopation, tied notes, or an upbeat that students misread as a downbeat, causing entries to fall one beat late. Incorrectly interpreted articulation markings (e.g., staccato vs. tenuto) or phrasing slurs could cause uniform articulation errors. \newline - Cueing and ensemble leadership: The section may be waiting for a visual or aural cue from the conductor or another section and misinterpreting the cue timing (e.g., following the end of a previous phrase rather than the intended beat). If a section leader cues late, the rest will follow late. \newline - Tempo perception and subdivision: Students may be hearing larger beats and not subdividing correctly; a fast tempo with weak internal subdivision can push their entries to the next beat. If they're counting in a way that emphasizes the wrong pulse, entries slip by one beat. \newline \textit{[\ldots\ truncated \ldots]} \newline In all three suggestions, provide immediate feedback language the instructor can use: concise, prescriptive prompts such as "Watch my prep---one-and-two-and---play on one,Slow---exaggerate staccato---three times," or "Principal breath---everyone in sync---play." These brief commands reduce ambiguity and cognitive load. \newline Final check and next steps: After implementing one or two of the corrections, the instructor should run the passage twice at tempo and listen for consistent on-beat attacks and uniform articulation. If timing errors persist but articulation improves, focus next on tempo/subdivision drills. If timing clears but articulation remains, do isolated articulation exercises (e.g., repeated-note staccato at metronome). Document whether the cause appears primarily notational misunderstanding, cueing, tempo perception, or technique so the instructor can schedule follow-up micro-rehearsals addressing the root cause. \\
  \midrule
  \textbf{Kimi-K2.6} & score 4/10 \\
  & [3-4] The error analysis is generally accurate and well-structured, but the answer is incomplete and cuts off mid-sentence. Crucially, it entirely omits the required three concrete corrective suggestions, implementation steps, and measurable cues, making the response inapplicable for direct use in the lesson context. \\
  \textbf{Innospark-1T} & score 7/10 \\
  & Let me think through this carefully. \newline The question/task: You are a teaching assistant monitoring a university-level ensemble rehearsal. During a live performance, a section of student musicians consistently plays a passage one beat late and with incorrect articulation compared to the score. Describe, in detail, how you would detect and analyze the root causes of this error in real ti \newline My reasoning: \textless{}Thinking\textgreater{} \newline Let me work through this systematically: \newline 1. Initial Detection \& Analysis Framework \newline I need to observe the specific error pattern: one beat late + incorrect articulation in a section. Score check: Are they reading correctly? Root Cause Analysis Framework \newline \textit{[\ldots\ truncated \ldots]} \newline | Notation | Misread articulation | Examine student scores | Explicit marking instruction | \newline This systematic approach allows for rapid diagnosis and targeted intervention within the rehearsal time constraints. \\
  \specialrule{1.2pt}{8pt}{8pt}
  \bottomrule
\end{longtable}